% =========================
% CHAPITRE 3 : CONCEPTION
% Fichier a inclure depuis le rapport principal via % =========================
% CHAPITRE 3 : CONCEPTION
% Fichier a inclure depuis le rapport principal via % =========================
% CHAPITRE 3 : CONCEPTION
% Fichier a inclure depuis le rapport principal via % =========================
% CHAPITRE 3 : CONCEPTION
% Fichier a inclure depuis le rapport principal via \input{chapters/chapitre3_conception}
% =========================

\chapter{Conception}
\label{chap:conception}

\section{Introduction}
Ce chapitre presente la conception de l'application de depistage automatise de la retinopathie diabetique.
Il transforme les besoins metier en une specification technique claire, en passant par l'analyse du systeme,
la modelisation des exigences et la definition de l'architecture.

L'objectif est de garantir une solution coherent entre la pratique clinique et les contraintes informatiques,
avec une separation nette des responsabilites entre les differents profils utilisateurs.

\section{Langage et outils de modelisation}

\subsection{Langage de modelisation}
Le langage retenu pour la conception est UML (Unified Modeling Language), car il permet de representer
le systeme selon plusieurs vues complementaires :
\begin{itemize}
    \item vue fonctionnelle (diagrammes de cas d'utilisation) ;
    \item vue statique (diagramme de classes) ;
    \item vue dynamique (diagrammes de sequence).
\end{itemize}

UML est adapte a ce projet car il facilite la communication entre les acteurs techniques et metier,
et fournit une base de reference pour l'implementation.

\subsection{Outils de modelisation}
Les outils utilises pour produire et maintenir les diagrammes sont :
\begin{itemize}
    \item Mermaid (fichiers .mmd) pour versionner les diagrammes avec le code source ;
    \item draw.io / diagrams.net pour les exports visuels a integrer dans le memoire ;
    \item LaTeX pour la redaction finale et la tracabilite des sections.
\end{itemize}

Le choix de Mermaid permet une meilleure reproductibilite des diagrammes et une maintenance plus simple
au cours des iterations du projet.

\section{Analyse de systeme}

\subsection{Acteurs}
L'analyse metier identifie des acteurs principaux (metier) et des acteurs secondaires (techniques)
qui interagissent avec l'application :
\begin{table}[H]
\centering
\caption{Acteurs du systeme et principales taches}
\label{tab:acteurs_systeme}
\begin{tabular}{|p{2.8cm}|p{3.2cm}|p{8.3cm}|}
\hline
\textbf{Acteur} & \textbf{Type d'acteur} & \textbf{Taches} \\
\hline
\textbf{CenterAdmin} & Acteur metier principal &
\begin{tabular}[t]{@{}l@{}}
- Creer, consulter, modifier et archiver les fiches patient \\
- Gerer les comptes medecins rattaches au centre \\
- Verifier la qualite de l'image avant soumission \\
- Lancer l'examen et suivre son etat (en cours, termine) \\
- Consulter l'historique des examens du centre \\
- Coordonner avec le medecin en cas de resultat severe
\end{tabular} \\
\hline
\textbf{Doctor} & Acteur metier principal &
\begin{tabular}[t]{@{}l@{}}
- Consulter les examens classes par priorite clinique \\
- Analyser le grade IA, le score de confiance et la heatmap Grad-CAM \\
- Interpreter les zones suspectes avec le contexte medical du patient \\
- Valider la conduite clinique (surveillance, controle rapide, orientation) \\
- Traiter les alertes urgentes et les marquer comme resolues \\
- Ajouter des notes cliniques pour la tracabilite
\end{tabular} \\
\hline
\textbf{SuperAdmin} & Acteur de gouvernance &
\begin{tabular}[t]{@{}l@{}}
- Administrer les centres (activation, desactivation, parametrage) \\
- Gerer les comptes globaux et les droits d'acces \\
- Superviser les indicateurs globaux (volumetrie, alertes, disponibilite) \\
- Suivre les tickets de support et les incidents critiques \\
- Coordonner les actions correctives avec les equipes techniques \\
- Assurer la conformite et la securite operationnelle
\end{tabular} \\
\hline
\textbf{Service IA} & Acteur technique secondaire &
\begin{tabular}[t]{@{}l@{}}
- Recevoir l'image du fond d'oeil depuis le backend \\
- Executer l'inference pour predire le grade de retinopathie (0--4) \\
- Calculer et retourner le niveau de confiance \\
- Generer la heatmap Grad-CAM pour l'explicabilite \\
- Renvoyer une reponse structuree exploitable par l'application
\end{tabular} \\
\hline
\textbf{Service WebSocket} & Acteur technique secondaire &
\begin{tabular}[t]{@{}l@{}}
- Maintenir les connexions temps reel avec les interfaces clientes \\
- Authentifier les connexions autorisees \\
- Diffuser les notifications de nouveaux examens/alertes \\
- Assurer la mise a jour immediate des tableaux de bord
\end{tabular} \\
\hline
\textbf{Serveur SMTP} & Acteur technique secondaire &
\begin{tabular}[t]{@{}l@{}}
- Recevoir les demandes d'envoi d'e-mails depuis le backend \\
- Envoyer les notifications pour les cas severes \\
- Fournir un canal de communication asynchrone en cas d'urgence
\end{tabular} \\
\hline
\end{tabular}
\end{table}

\textbf{Important :} \textit{Patient} et \textit{Center} ne sont pas modelises comme acteurs,
mais comme entites metier manipulees par le systeme.

\subsection{Exigences fonctionnelles}
Les exigences fonctionnelles sont regroupees par domaine :

\subsubsection*{Gestion des comptes et securite}
\begin{itemize}
    \item authentification par role (CenterAdmin, Doctor, SuperAdmin) ;
    \item gestion de session et controle d'acces selon le profil ;
    \item deconnexion securisee.
\end{itemize}

\subsubsection*{Gestion clinique et operationnelle}
\begin{itemize}
    \item creation, consultation, modification et archivage des patients ;
    \item gestion des medecins par centre ;
    \item creation d'un examen avec upload de l'image du fond d'oeil.
\end{itemize}

\subsubsection*{Analyse IA et explicabilite}
\begin{itemize}
    \item prediction automatique du grade de retinopathie (0 a 4) ;
    \item calcul d'un score de confiance associe a la prediction ;
    \item generation d'une carte de chaleur Grad-CAM ;
    \item conservation de la sortie IA dans l'historique d'examens.
\end{itemize}

\subsubsection*{Alerte et supervision}
\begin{itemize}
    \item creation d'alertes automatiques pour les cas severes (grade >= 3) ;
    \item notification temps reel du medecin ;
    \item suivi des alertes (non lue, lue, resolue) ;
    \item supervision globale et support par le SuperAdmin.
\end{itemize}

\subsection{Exigences non fonctionnelles}
Les exigences non fonctionnelles assurent la qualite globale du systeme :
\begin{itemize}
    \item \textbf{Performance} : temps de reponse court pour l'inference et l'affichage des resultats ;
    \item \textbf{Disponibilite} : acces continu a l'application pour les centres et medecins ;
    \item \textbf{Securite} : protection des donnees medicales, isolation par centre, controle des acces ;
    \item \textbf{Fiabilite} : tracabilite des examens et des alertes ;
    \item \textbf{Scalabilite} : possibilite d'ajouter des centres et d'augmenter le volume d'examens ;
    \item \textbf{Maintenabilite} : architecture modulaire (web app, service IA, service temps reel) ;
    \item \textbf{Explicabilite} : utilisation d'une visualisation interpretable pour appuyer la decision medicale.
\end{itemize}

Dans la configuration courante du projet, la visualisation explicative par defaut est une Grad-CAM
focalisee sur la couche \textit{layer4} du modele ResNet, afin de produire des cartes de chaleur plus ciblees.

\section{Modelisation des exigences fonctionnelles}

\subsection{Diagramme de cas d'utilisation general}
Le diagramme general regroupe les interactions majeures entre les trois acteurs et le systeme.

\textbf{Cas d'utilisation communs :}
\begin{itemize}
    \item se connecter ;
    \item consulter son tableau de bord ;
    \item se deconnecter.
\end{itemize}

\textbf{Cas d'utilisation CenterAdmin :}
\begin{itemize}
    \item gerer patients ;
    \item gerer medecins ;
    \item creer examen ;
    \item consulter historique des examens du centre.
\end{itemize}

\textbf{Cas d'utilisation Doctor :}
\begin{itemize}
    \item consulter liste des examens ;
    \item consulter detail d'un examen ;
    \item visualiser prediction et heatmap ;
    \item traiter alertes cliniques.
\end{itemize}

\textbf{Cas d'utilisation SuperAdmin :}
\begin{itemize}
    \item gerer centres ;
    \item gerer comptes globaux ;
    \item superviser l'activite de la plateforme ;
    \item gerer les tickets de support et le suivi global.
\end{itemize}

\subsection{Diagrammes de cas d'utilisation detaillee}

\subsubsection*{a) Use case Admin (CenterAdmin)}
Le profil CenterAdmin assure la preparation operationnelle du depistage.
Le scenario principal est le suivant :
\begin{enumerate}
    \item l'admin authentifie ouvre l'espace centre ;
    \item il cree ou selectionne un patient ;
    \item il associe un medecin referent ;
    \item il charge l'image retinienne ;
    \item il valide la creation d'examen ;
    \item le systeme declenche l'analyse IA et enregistre le resultat.
\end{enumerate}

\begin{figure}[H]
\centering
\begin{tikzpicture}[scale=0.90, every node/.style={font=\small}]
    % Limite du systeme
    \draw[thick] (1.7,0.8) rectangle (13.3,8.3);
    \node[anchor=west] at (1.9,8.02) {Systeme de depistage RD};

    % Acteur CenterAdmin
    \draw (0.1,6.0) circle (0.3);
    \draw (0.1,5.7) -- (0.1,4.8);
    \draw (-0.45,5.35) -- (0.65,5.35);
    \draw (0.1,4.8) -- (-0.35,4.1);
    \draw (0.1,4.8) -- (0.55,4.1);
    \node at (0.1,3.7) {CenterAdmin};

    % Acteur secondaire Service IA
    \draw (14.8,5.0) circle (0.3);
    \draw (14.8,4.7) -- (14.8,3.8);
    \draw (14.25,4.35) -- (15.35,4.35);
    \draw (14.8,3.8) -- (14.35,3.1);
    \draw (14.8,3.8) -- (15.25,3.1);
    \node[align=center] at (14.8,2.55) {Service IA};

    % Cas d'utilisation principaux
    \node[draw, ellipse, minimum width=4.8cm, minimum height=0.9cm] (ua1) at (7.6,7.1) {Se connecter};
    \node[draw, ellipse, minimum width=4.8cm, minimum height=0.9cm] (ua2) at (7.6,6.1) {Gerer patients};
    \node[draw, ellipse, minimum width=4.8cm, minimum height=0.9cm] (ua3) at (7.6,5.1) {Gerer medecins};
    \node[draw, ellipse, minimum width=4.8cm, minimum height=0.9cm] (ua4) at (7.6,4.1) {Creer examen};
    \node[draw, ellipse, minimum width=4.8cm, minimum height=0.9cm] (ua5) at (7.6,3.1) {Consulter historique des examens};
    \node[draw, ellipse, minimum width=4.8cm, minimum height=0.9cm] (ua6) at (7.6,2.1) {Se deconnecter};

    % Sous-cas d'utilisation de "Creer examen"
    \node[draw, ellipse, minimum width=4.2cm, minimum height=0.85cm] (ua41) at (4.2,1.25) {Televerser image retinienne};
    \node[draw, ellipse, minimum width=4.2cm, minimum height=0.85cm] (ua42) at (7.6,1.25) {Associer patient et medecin};
    \node[draw, ellipse, minimum width=4.2cm, minimum height=0.85cm] (ua43) at (11.0,1.25) {Soumettre a l'analyse IA};

    % Liens acteur principal
    \draw (0.65,5.35) -- (ua1.west);
    \draw (0.65,5.35) -- (ua2.west);
    \draw (0.65,5.35) -- (ua3.west);
    \draw (0.65,5.35) -- (ua4.west);
    \draw (0.65,5.35) -- (ua5.west);
    \draw (0.65,5.35) -- (ua6.west);

    % Include
    \draw[->, dashed] (ua4.south west) -- (ua41.north);
    \node at (5.6,2.4) {<<include>>};
    \draw[->, dashed] (ua4.south) -- (ua42.north);
    \node at (7.6,2.4) {<<include>>};
    \draw[->, dashed] (ua4.south east) -- (ua43.north);
    \node at (9.7,2.4) {<<include>>};

    % Lien avec l'acteur secondaire IA
    \draw (14.25,4.35) -- (ua43.east);
\end{tikzpicture}
\caption{Diagramme de cas d'utilisation detaille -- acteur CenterAdmin.}
\label{fig:usecase_centeradmin}
\end{figure}

\textbf{Post-condition :} un examen complet (image, grade, confiance, heatmap) est disponible.

\subsubsection*{b) Use case Doctor}
Le profil Doctor exploite les resultats pour l'aide a la decision clinique.
Le scenario principal est :
\begin{enumerate}
    \item le medecin consulte la liste des examens classes par severite ;
    \item il ouvre un examen ;
    \item il analyse le grade predit, la confiance et la heatmap ;
    \item il decide la conduite clinique (surveillance, controle rapide, orientation urgente) ;
    \item il marque l'alerte comme traitee si necessaire.
\end{enumerate}

\begin{figure}[H]
\centering
\begin{tikzpicture}[scale=0.92, every node/.style={font=\small}]
    % Limite du systeme
    \draw[thick] (1.5,0.8) rectangle (12.5,8.2);
    \node[anchor=west] at (1.7,7.95) {Systeme de depistage RD};

    % Acteur Doctor
    \draw (0,5.9) circle (0.3);
    \draw (0,5.6) -- (0,4.7);
    \draw (-0.55,5.25) -- (0.55,5.25);
    \draw (0,4.7) -- (-0.45,4.0);
    \draw (0,4.7) -- (0.45,4.0);
    \node at (0,3.6) {Doctor};

    % Cas d'utilisation
    \node[draw, ellipse, minimum width=4.9cm, minimum height=0.95cm] (uc1) at (7.0,6.9) {Se connecter};
    \node[draw, ellipse, minimum width=4.9cm, minimum height=0.95cm] (uc2) at (7.0,5.8) {Consulter liste des examens};
    \node[draw, ellipse, minimum width=4.9cm, minimum height=0.95cm] (uc3) at (7.0,4.7) {Consulter detail d'un examen};
    \node[draw, ellipse, minimum width=4.9cm, minimum height=0.95cm] (uc4) at (7.0,3.6) {Visualiser prediction et heatmap};
    \node[draw, ellipse, minimum width=4.9cm, minimum height=0.95cm] (uc5) at (7.0,2.5) {Traiter alertes cliniques};
    \node[draw, ellipse, minimum width=4.9cm, minimum height=0.95cm] (uc6) at (7.0,1.4) {Se deconnecter};

    % Associations acteur -> use cases
    \draw (0.55,5.25) -- (uc1.west);
    \draw (0.55,5.25) -- (uc2.west);
    \draw (0.55,5.25) -- (uc3.west);
    \draw (0.55,5.25) -- (uc4.west);
    \draw (0.55,5.25) -- (uc5.west);
    \draw (0.55,5.25) -- (uc6.west);

    % Relations include
    \draw[->, dashed] (uc3.south) -- (uc4.north);
    \node at (10.8,4.15) {<<include>>};
    \draw[->, dashed] (uc5.north) -- (uc3.south);
    \node at (10.8,3.05) {<<include>>};
\end{tikzpicture}
\caption{Diagramme de cas d'utilisation detaille -- acteur Doctor.}
\label{fig:usecase_doctor}
\end{figure}

\textbf{Post-condition :} la decision medicale est tracee et l'etat d'alerte est mis a jour.

\subsubsection*{c) Use case SuperAdmin}
Le profil SuperAdmin intervient au niveau de la gouvernance de la plateforme.
Le scenario principal est :
\begin{enumerate}
    \item le super administrateur se connecte au dashboard global ;
    \item il supervise les centres actifs et les comptes utilisateurs ;
    \item il controle la qualite de service (tickets, incidents, disponibilite) ;
    \item il applique les actions d'administration necessaires.
\end{enumerate}

\textbf{Post-condition :} la plateforme reste operationnelle, securisee et conforme aux objectifs de service.

\subsection{Diagramme de classes}
Le diagramme de classes represente les entites de base et leurs relations.
Les classes principales sont :
\begin{itemize}
    \item \textbf{Center} ;
    \item \textbf{User} (specialise en CenterAdmin, Doctor, SuperAdmin selon le contexte applicatif) ;
    \item \textbf{Patient} ;
    \item \textbf{Exam} ;
    \item \textbf{Alert} ;
    \item \textbf{SupportTicket} (niveau supervision).
\end{itemize}

Relations essentielles :
\begin{itemize}
    \item un Center possede plusieurs Users et plusieurs Patients ;
    \item un Patient possede plusieurs Exams ;
    \item un Exam est cree par un User et peut generer zero ou plusieurs Alerts ;
    \item une Alert est assignee a un Doctor pour traitement ;
    \item un SupportTicket est suivi par le SuperAdmin.
\end{itemize}

Cette modelisation clarifie le cycle de vie des donnees medicales, depuis l'acquisition de l'image
jusqu'au traitement des alertes.

\section{Architecture complete du systeme}

\subsection{Architecture globale}
L'architecture adoptee est une architecture modulaire multi-composants :
\begin{itemize}
    \item \textbf{Frontend Web} : interfaces de connexion, gestion metier et visualisation ;
    \item \textbf{Backend applicatif} : logique metier, API, securite, persistence ;
    \item \textbf{Service IA} : inference de grade RD et generation Grad-CAM ;
    \item \textbf{Service temps reel} : diffusion des notifications et alertes ;
    \item \textbf{Base de donnees} : stockage des utilisateurs, patients, examens et alertes.
\end{itemize}

Ce decouplage facilite la maintenance, les tests et l'evolution independante des composants.

\subsection{Architecture detaillee (diagrammes de sequence)}
Les diagrammes de sequence detaillent les echanges inter-services pour les cas critiques.

\subsubsection*{Sequence 1 : creation d'un examen}
\begin{enumerate}
    \item CenterAdmin envoie l'image au Backend ;
    \item le Backend stocke l'image et appelle le Service IA ;
    \item le Service IA retourne grade, confiance et heatmap ;
    \item le Backend enregistre l'examen ;
    \item si grade >= 3, le Backend cree une alerte ;
    \item le Service temps reel notifie le Doctor ;
    \item l'interface met a jour l'etat de maniere immediate.
\end{enumerate}

\subsubsection*{Sequence 2 : consultation et traitement d'alerte}
\begin{enumerate}
    \item le Doctor ouvre la liste des alertes ;
    \item le Backend retourne les alertes actives ;
    \item le Doctor consulte le detail d'examen ;
    \item il valide l'action clinique ;
    \item l'alerte passe a l'etat resolu avec historisation.
\end{enumerate}

\subsubsection*{Sequence 3 : supervision plateforme}
\begin{enumerate}
    \item le SuperAdmin consulte les indicateurs globaux ;
    \item le systeme agrege les statistiques multi-centres ;
    \item le SuperAdmin traite les tickets critiques et incidents ;
    \item les actions de supervision sont journalisees.
\end{enumerate}

\section{Conclusion}
Ce chapitre a presente la conception fonctionnelle et technique de la solution.
L'analyse des acteurs, la specification des exigences, les diagrammes UML et l'architecture complete
fournissent un cadre clair pour l'implementation.

La phase suivante peut s'appuyer sur cette conception pour la realisation detaillee,
les tests d'integration et la validation finale de la plateforme en contexte reel.

% =========================

\chapter{Conception}
\label{chap:conception}

\section{Introduction}
Ce chapitre presente la conception de l'application de depistage automatise de la retinopathie diabetique.
Il transforme les besoins metier en une specification technique claire, en passant par l'analyse du systeme,
la modelisation des exigences et la definition de l'architecture.

L'objectif est de garantir une solution coherent entre la pratique clinique et les contraintes informatiques,
avec une separation nette des responsabilites entre les differents profils utilisateurs.

\section{Langage et outils de modelisation}

\subsection{Langage de modelisation}
Le langage retenu pour la conception est UML (Unified Modeling Language), car il permet de representer
le systeme selon plusieurs vues complementaires :
\begin{itemize}
    \item vue fonctionnelle (diagrammes de cas d'utilisation) ;
    \item vue statique (diagramme de classes) ;
    \item vue dynamique (diagrammes de sequence).
\end{itemize}

UML est adapte a ce projet car il facilite la communication entre les acteurs techniques et metier,
et fournit une base de reference pour l'implementation.

\subsection{Outils de modelisation}
Les outils utilises pour produire et maintenir les diagrammes sont :
\begin{itemize}
    \item Mermaid (fichiers .mmd) pour versionner les diagrammes avec le code source ;
    \item draw.io / diagrams.net pour les exports visuels a integrer dans le memoire ;
    \item LaTeX pour la redaction finale et la tracabilite des sections.
\end{itemize}

Le choix de Mermaid permet une meilleure reproductibilite des diagrammes et une maintenance plus simple
au cours des iterations du projet.

\section{Analyse de systeme}

\subsection{Acteurs}
L'analyse metier identifie des acteurs principaux (metier) et des acteurs secondaires (techniques)
qui interagissent avec l'application :
\begin{table}[H]
\centering
\caption{Acteurs du systeme et principales taches}
\label{tab:acteurs_systeme}
\begin{tabular}{|p{2.8cm}|p{3.2cm}|p{8.3cm}|}
\hline
\textbf{Acteur} & \textbf{Type d'acteur} & \textbf{Taches} \\
\hline
\textbf{CenterAdmin} & Acteur metier principal &
\begin{tabular}[t]{@{}l@{}}
- Creer, consulter, modifier et archiver les fiches patient \\
- Gerer les comptes medecins rattaches au centre \\
- Verifier la qualite de l'image avant soumission \\
- Lancer l'examen et suivre son etat (en cours, termine) \\
- Consulter l'historique des examens du centre \\
- Coordonner avec le medecin en cas de resultat severe
\end{tabular} \\
\hline
\textbf{Doctor} & Acteur metier principal &
\begin{tabular}[t]{@{}l@{}}
- Consulter les examens classes par priorite clinique \\
- Analyser le grade IA, le score de confiance et la heatmap Grad-CAM \\
- Interpreter les zones suspectes avec le contexte medical du patient \\
- Valider la conduite clinique (surveillance, controle rapide, orientation) \\
- Traiter les alertes urgentes et les marquer comme resolues \\
- Ajouter des notes cliniques pour la tracabilite
\end{tabular} \\
\hline
\textbf{SuperAdmin} & Acteur de gouvernance &
\begin{tabular}[t]{@{}l@{}}
- Administrer les centres (activation, desactivation, parametrage) \\
- Gerer les comptes globaux et les droits d'acces \\
- Superviser les indicateurs globaux (volumetrie, alertes, disponibilite) \\
- Suivre les tickets de support et les incidents critiques \\
- Coordonner les actions correctives avec les equipes techniques \\
- Assurer la conformite et la securite operationnelle
\end{tabular} \\
\hline
\textbf{Service IA} & Acteur technique secondaire &
\begin{tabular}[t]{@{}l@{}}
- Recevoir l'image du fond d'oeil depuis le backend \\
- Executer l'inference pour predire le grade de retinopathie (0--4) \\
- Calculer et retourner le niveau de confiance \\
- Generer la heatmap Grad-CAM pour l'explicabilite \\
- Renvoyer une reponse structuree exploitable par l'application
\end{tabular} \\
\hline
\textbf{Service WebSocket} & Acteur technique secondaire &
\begin{tabular}[t]{@{}l@{}}
- Maintenir les connexions temps reel avec les interfaces clientes \\
- Authentifier les connexions autorisees \\
- Diffuser les notifications de nouveaux examens/alertes \\
- Assurer la mise a jour immediate des tableaux de bord
\end{tabular} \\
\hline
\textbf{Serveur SMTP} & Acteur technique secondaire &
\begin{tabular}[t]{@{}l@{}}
- Recevoir les demandes d'envoi d'e-mails depuis le backend \\
- Envoyer les notifications pour les cas severes \\
- Fournir un canal de communication asynchrone en cas d'urgence
\end{tabular} \\
\hline
\end{tabular}
\end{table}

\textbf{Important :} \textit{Patient} et \textit{Center} ne sont pas modelises comme acteurs,
mais comme entites metier manipulees par le systeme.

\subsection{Exigences fonctionnelles}
Les exigences fonctionnelles sont regroupees par domaine :

\subsubsection*{Gestion des comptes et securite}
\begin{itemize}
    \item authentification par role (CenterAdmin, Doctor, SuperAdmin) ;
    \item gestion de session et controle d'acces selon le profil ;
    \item deconnexion securisee.
\end{itemize}

\subsubsection*{Gestion clinique et operationnelle}
\begin{itemize}
    \item creation, consultation, modification et archivage des patients ;
    \item gestion des medecins par centre ;
    \item creation d'un examen avec upload de l'image du fond d'oeil.
\end{itemize}

\subsubsection*{Analyse IA et explicabilite}
\begin{itemize}
    \item prediction automatique du grade de retinopathie (0 a 4) ;
    \item calcul d'un score de confiance associe a la prediction ;
    \item generation d'une carte de chaleur Grad-CAM ;
    \item conservation de la sortie IA dans l'historique d'examens.
\end{itemize}

\subsubsection*{Alerte et supervision}
\begin{itemize}
    \item creation d'alertes automatiques pour les cas severes (grade >= 3) ;
    \item notification temps reel du medecin ;
    \item suivi des alertes (non lue, lue, resolue) ;
    \item supervision globale et support par le SuperAdmin.
\end{itemize}

\subsection{Exigences non fonctionnelles}
Les exigences non fonctionnelles assurent la qualite globale du systeme :
\begin{itemize}
    \item \textbf{Performance} : temps de reponse court pour l'inference et l'affichage des resultats ;
    \item \textbf{Disponibilite} : acces continu a l'application pour les centres et medecins ;
    \item \textbf{Securite} : protection des donnees medicales, isolation par centre, controle des acces ;
    \item \textbf{Fiabilite} : tracabilite des examens et des alertes ;
    \item \textbf{Scalabilite} : possibilite d'ajouter des centres et d'augmenter le volume d'examens ;
    \item \textbf{Maintenabilite} : architecture modulaire (web app, service IA, service temps reel) ;
    \item \textbf{Explicabilite} : utilisation d'une visualisation interpretable pour appuyer la decision medicale.
\end{itemize}

Dans la configuration courante du projet, la visualisation explicative par defaut est une Grad-CAM
focalisee sur la couche \textit{layer4} du modele ResNet, afin de produire des cartes de chaleur plus ciblees.

\section{Modelisation des exigences fonctionnelles}

\subsection{Diagramme de cas d'utilisation general}
Le diagramme general regroupe les interactions majeures entre les trois acteurs et le systeme.

\textbf{Cas d'utilisation communs :}
\begin{itemize}
    \item se connecter ;
    \item consulter son tableau de bord ;
    \item se deconnecter.
\end{itemize}

\textbf{Cas d'utilisation CenterAdmin :}
\begin{itemize}
    \item gerer patients ;
    \item gerer medecins ;
    \item creer examen ;
    \item consulter historique des examens du centre.
\end{itemize}

\textbf{Cas d'utilisation Doctor :}
\begin{itemize}
    \item consulter liste des examens ;
    \item consulter detail d'un examen ;
    \item visualiser prediction et heatmap ;
    \item traiter alertes cliniques.
\end{itemize}

\textbf{Cas d'utilisation SuperAdmin :}
\begin{itemize}
    \item gerer centres ;
    \item gerer comptes globaux ;
    \item superviser l'activite de la plateforme ;
    \item gerer les tickets de support et le suivi global.
\end{itemize}

\subsection{Diagrammes de cas d'utilisation detaillee}

\subsubsection*{a) Use case Admin (CenterAdmin)}
Le profil CenterAdmin assure la preparation operationnelle du depistage.
Le scenario principal est le suivant :
\begin{enumerate}
    \item l'admin authentifie ouvre l'espace centre ;
    \item il cree ou selectionne un patient ;
    \item il associe un medecin referent ;
    \item il charge l'image retinienne ;
    \item il valide la creation d'examen ;
    \item le systeme declenche l'analyse IA et enregistre le resultat.
\end{enumerate}

\begin{figure}[H]
\centering
\begin{tikzpicture}[scale=0.90, every node/.style={font=\small}]
    % Limite du systeme
    \draw[thick] (1.7,0.8) rectangle (13.3,8.3);
    \node[anchor=west] at (1.9,8.02) {Systeme de depistage RD};

    % Acteur CenterAdmin
    \draw (0.1,6.0) circle (0.3);
    \draw (0.1,5.7) -- (0.1,4.8);
    \draw (-0.45,5.35) -- (0.65,5.35);
    \draw (0.1,4.8) -- (-0.35,4.1);
    \draw (0.1,4.8) -- (0.55,4.1);
    \node at (0.1,3.7) {CenterAdmin};

    % Acteur secondaire Service IA
    \draw (14.8,5.0) circle (0.3);
    \draw (14.8,4.7) -- (14.8,3.8);
    \draw (14.25,4.35) -- (15.35,4.35);
    \draw (14.8,3.8) -- (14.35,3.1);
    \draw (14.8,3.8) -- (15.25,3.1);
    \node[align=center] at (14.8,2.55) {Service IA};

    % Cas d'utilisation principaux
    \node[draw, ellipse, minimum width=4.8cm, minimum height=0.9cm] (ua1) at (7.6,7.1) {Se connecter};
    \node[draw, ellipse, minimum width=4.8cm, minimum height=0.9cm] (ua2) at (7.6,6.1) {Gerer patients};
    \node[draw, ellipse, minimum width=4.8cm, minimum height=0.9cm] (ua3) at (7.6,5.1) {Gerer medecins};
    \node[draw, ellipse, minimum width=4.8cm, minimum height=0.9cm] (ua4) at (7.6,4.1) {Creer examen};
    \node[draw, ellipse, minimum width=4.8cm, minimum height=0.9cm] (ua5) at (7.6,3.1) {Consulter historique des examens};
    \node[draw, ellipse, minimum width=4.8cm, minimum height=0.9cm] (ua6) at (7.6,2.1) {Se deconnecter};

    % Sous-cas d'utilisation de "Creer examen"
    \node[draw, ellipse, minimum width=4.2cm, minimum height=0.85cm] (ua41) at (4.2,1.25) {Televerser image retinienne};
    \node[draw, ellipse, minimum width=4.2cm, minimum height=0.85cm] (ua42) at (7.6,1.25) {Associer patient et medecin};
    \node[draw, ellipse, minimum width=4.2cm, minimum height=0.85cm] (ua43) at (11.0,1.25) {Soumettre a l'analyse IA};

    % Liens acteur principal
    \draw (0.65,5.35) -- (ua1.west);
    \draw (0.65,5.35) -- (ua2.west);
    \draw (0.65,5.35) -- (ua3.west);
    \draw (0.65,5.35) -- (ua4.west);
    \draw (0.65,5.35) -- (ua5.west);
    \draw (0.65,5.35) -- (ua6.west);

    % Include
    \draw[->, dashed] (ua4.south west) -- (ua41.north);
    \node at (5.6,2.4) {<<include>>};
    \draw[->, dashed] (ua4.south) -- (ua42.north);
    \node at (7.6,2.4) {<<include>>};
    \draw[->, dashed] (ua4.south east) -- (ua43.north);
    \node at (9.7,2.4) {<<include>>};

    % Lien avec l'acteur secondaire IA
    \draw (14.25,4.35) -- (ua43.east);
\end{tikzpicture}
\caption{Diagramme de cas d'utilisation detaille -- acteur CenterAdmin.}
\label{fig:usecase_centeradmin}
\end{figure}

\textbf{Post-condition :} un examen complet (image, grade, confiance, heatmap) est disponible.

\subsubsection*{b) Use case Doctor}
Le profil Doctor exploite les resultats pour l'aide a la decision clinique.
Le scenario principal est :
\begin{enumerate}
    \item le medecin consulte la liste des examens classes par severite ;
    \item il ouvre un examen ;
    \item il analyse le grade predit, la confiance et la heatmap ;
    \item il decide la conduite clinique (surveillance, controle rapide, orientation urgente) ;
    \item il marque l'alerte comme traitee si necessaire.
\end{enumerate}

\begin{figure}[H]
\centering
\begin{tikzpicture}[scale=0.92, every node/.style={font=\small}]
    % Limite du systeme
    \draw[thick] (1.5,0.8) rectangle (12.5,8.2);
    \node[anchor=west] at (1.7,7.95) {Systeme de depistage RD};

    % Acteur Doctor
    \draw (0,5.9) circle (0.3);
    \draw (0,5.6) -- (0,4.7);
    \draw (-0.55,5.25) -- (0.55,5.25);
    \draw (0,4.7) -- (-0.45,4.0);
    \draw (0,4.7) -- (0.45,4.0);
    \node at (0,3.6) {Doctor};

    % Cas d'utilisation
    \node[draw, ellipse, minimum width=4.9cm, minimum height=0.95cm] (uc1) at (7.0,6.9) {Se connecter};
    \node[draw, ellipse, minimum width=4.9cm, minimum height=0.95cm] (uc2) at (7.0,5.8) {Consulter liste des examens};
    \node[draw, ellipse, minimum width=4.9cm, minimum height=0.95cm] (uc3) at (7.0,4.7) {Consulter detail d'un examen};
    \node[draw, ellipse, minimum width=4.9cm, minimum height=0.95cm] (uc4) at (7.0,3.6) {Visualiser prediction et heatmap};
    \node[draw, ellipse, minimum width=4.9cm, minimum height=0.95cm] (uc5) at (7.0,2.5) {Traiter alertes cliniques};
    \node[draw, ellipse, minimum width=4.9cm, minimum height=0.95cm] (uc6) at (7.0,1.4) {Se deconnecter};

    % Associations acteur -> use cases
    \draw (0.55,5.25) -- (uc1.west);
    \draw (0.55,5.25) -- (uc2.west);
    \draw (0.55,5.25) -- (uc3.west);
    \draw (0.55,5.25) -- (uc4.west);
    \draw (0.55,5.25) -- (uc5.west);
    \draw (0.55,5.25) -- (uc6.west);

    % Relations include
    \draw[->, dashed] (uc3.south) -- (uc4.north);
    \node at (10.8,4.15) {<<include>>};
    \draw[->, dashed] (uc5.north) -- (uc3.south);
    \node at (10.8,3.05) {<<include>>};
\end{tikzpicture}
\caption{Diagramme de cas d'utilisation detaille -- acteur Doctor.}
\label{fig:usecase_doctor}
\end{figure}

\textbf{Post-condition :} la decision medicale est tracee et l'etat d'alerte est mis a jour.

\subsubsection*{c) Use case SuperAdmin}
Le profil SuperAdmin intervient au niveau de la gouvernance de la plateforme.
Le scenario principal est :
\begin{enumerate}
    \item le super administrateur se connecte au dashboard global ;
    \item il supervise les centres actifs et les comptes utilisateurs ;
    \item il controle la qualite de service (tickets, incidents, disponibilite) ;
    \item il applique les actions d'administration necessaires.
\end{enumerate}

\textbf{Post-condition :} la plateforme reste operationnelle, securisee et conforme aux objectifs de service.

\subsection{Diagramme de classes}
Le diagramme de classes represente les entites de base et leurs relations.
Les classes principales sont :
\begin{itemize}
    \item \textbf{Center} ;
    \item \textbf{User} (specialise en CenterAdmin, Doctor, SuperAdmin selon le contexte applicatif) ;
    \item \textbf{Patient} ;
    \item \textbf{Exam} ;
    \item \textbf{Alert} ;
    \item \textbf{SupportTicket} (niveau supervision).
\end{itemize}

Relations essentielles :
\begin{itemize}
    \item un Center possede plusieurs Users et plusieurs Patients ;
    \item un Patient possede plusieurs Exams ;
    \item un Exam est cree par un User et peut generer zero ou plusieurs Alerts ;
    \item une Alert est assignee a un Doctor pour traitement ;
    \item un SupportTicket est suivi par le SuperAdmin.
\end{itemize}

Cette modelisation clarifie le cycle de vie des donnees medicales, depuis l'acquisition de l'image
jusqu'au traitement des alertes.

\section{Architecture complete du systeme}

\subsection{Architecture globale}
L'architecture adoptee est une architecture modulaire multi-composants :
\begin{itemize}
    \item \textbf{Frontend Web} : interfaces de connexion, gestion metier et visualisation ;
    \item \textbf{Backend applicatif} : logique metier, API, securite, persistence ;
    \item \textbf{Service IA} : inference de grade RD et generation Grad-CAM ;
    \item \textbf{Service temps reel} : diffusion des notifications et alertes ;
    \item \textbf{Base de donnees} : stockage des utilisateurs, patients, examens et alertes.
\end{itemize}

Ce decouplage facilite la maintenance, les tests et l'evolution independante des composants.

\subsection{Architecture detaillee (diagrammes de sequence)}
Les diagrammes de sequence detaillent les echanges inter-services pour les cas critiques.

\subsubsection*{Sequence 1 : creation d'un examen}
\begin{enumerate}
    \item CenterAdmin envoie l'image au Backend ;
    \item le Backend stocke l'image et appelle le Service IA ;
    \item le Service IA retourne grade, confiance et heatmap ;
    \item le Backend enregistre l'examen ;
    \item si grade >= 3, le Backend cree une alerte ;
    \item le Service temps reel notifie le Doctor ;
    \item l'interface met a jour l'etat de maniere immediate.
\end{enumerate}

\subsubsection*{Sequence 2 : consultation et traitement d'alerte}
\begin{enumerate}
    \item le Doctor ouvre la liste des alertes ;
    \item le Backend retourne les alertes actives ;
    \item le Doctor consulte le detail d'examen ;
    \item il valide l'action clinique ;
    \item l'alerte passe a l'etat resolu avec historisation.
\end{enumerate}

\subsubsection*{Sequence 3 : supervision plateforme}
\begin{enumerate}
    \item le SuperAdmin consulte les indicateurs globaux ;
    \item le systeme agrege les statistiques multi-centres ;
    \item le SuperAdmin traite les tickets critiques et incidents ;
    \item les actions de supervision sont journalisees.
\end{enumerate}

\section{Conclusion}
Ce chapitre a presente la conception fonctionnelle et technique de la solution.
L'analyse des acteurs, la specification des exigences, les diagrammes UML et l'architecture complete
fournissent un cadre clair pour l'implementation.

La phase suivante peut s'appuyer sur cette conception pour la realisation detaillee,
les tests d'integration et la validation finale de la plateforme en contexte reel.

% =========================

\chapter{Conception}
\label{chap:conception}

\section{Introduction}
Ce chapitre presente la conception de l'application de depistage automatise de la retinopathie diabetique.
Il transforme les besoins metier en une specification technique claire, en passant par l'analyse du systeme,
la modelisation des exigences et la definition de l'architecture.

L'objectif est de garantir une solution coherent entre la pratique clinique et les contraintes informatiques,
avec une separation nette des responsabilites entre les differents profils utilisateurs.

\section{Langage et outils de modelisation}

\subsection{Langage de modelisation}
Le langage retenu pour la conception est UML (Unified Modeling Language), car il permet de representer
le systeme selon plusieurs vues complementaires :
\begin{itemize}
    \item vue fonctionnelle (diagrammes de cas d'utilisation) ;
    \item vue statique (diagramme de classes) ;
    \item vue dynamique (diagrammes de sequence).
\end{itemize}

UML est adapte a ce projet car il facilite la communication entre les acteurs techniques et metier,
et fournit une base de reference pour l'implementation.

\subsection{Outils de modelisation}
Les outils utilises pour produire et maintenir les diagrammes sont :
\begin{itemize}
    \item Mermaid (fichiers .mmd) pour versionner les diagrammes avec le code source ;
    \item draw.io / diagrams.net pour les exports visuels a integrer dans le memoire ;
    \item LaTeX pour la redaction finale et la tracabilite des sections.
\end{itemize}

Le choix de Mermaid permet une meilleure reproductibilite des diagrammes et une maintenance plus simple
au cours des iterations du projet.

\section{Analyse de systeme}

\subsection{Acteurs}
L'analyse metier identifie des acteurs principaux (metier) et des acteurs secondaires (techniques)
qui interagissent avec l'application :
\begin{table}[H]
\centering
\caption{Acteurs du systeme et principales taches}
\label{tab:acteurs_systeme}
\begin{tabular}{|p{2.8cm}|p{3.2cm}|p{8.3cm}|}
\hline
\textbf{Acteur} & \textbf{Type d'acteur} & \textbf{Taches} \\
\hline
\textbf{CenterAdmin} & Acteur metier principal &
\begin{tabular}[t]{@{}l@{}}
- Creer, consulter, modifier et archiver les fiches patient \\
- Gerer les comptes medecins rattaches au centre \\
- Verifier la qualite de l'image avant soumission \\
- Lancer l'examen et suivre son etat (en cours, termine) \\
- Consulter l'historique des examens du centre \\
- Coordonner avec le medecin en cas de resultat severe
\end{tabular} \\
\hline
\textbf{Doctor} & Acteur metier principal &
\begin{tabular}[t]{@{}l@{}}
- Consulter les examens classes par priorite clinique \\
- Analyser le grade IA, le score de confiance et la heatmap Grad-CAM \\
- Interpreter les zones suspectes avec le contexte medical du patient \\
- Valider la conduite clinique (surveillance, controle rapide, orientation) \\
- Traiter les alertes urgentes et les marquer comme resolues \\
- Ajouter des notes cliniques pour la tracabilite
\end{tabular} \\
\hline
\textbf{SuperAdmin} & Acteur de gouvernance &
\begin{tabular}[t]{@{}l@{}}
- Administrer les centres (activation, desactivation, parametrage) \\
- Gerer les comptes globaux et les droits d'acces \\
- Superviser les indicateurs globaux (volumetrie, alertes, disponibilite) \\
- Suivre les tickets de support et les incidents critiques \\
- Coordonner les actions correctives avec les equipes techniques \\
- Assurer la conformite et la securite operationnelle
\end{tabular} \\
\hline
\textbf{Service IA} & Acteur technique secondaire &
\begin{tabular}[t]{@{}l@{}}
- Recevoir l'image du fond d'oeil depuis le backend \\
- Executer l'inference pour predire le grade de retinopathie (0--4) \\
- Calculer et retourner le niveau de confiance \\
- Generer la heatmap Grad-CAM pour l'explicabilite \\
- Renvoyer une reponse structuree exploitable par l'application
\end{tabular} \\
\hline
\textbf{Service WebSocket} & Acteur technique secondaire &
\begin{tabular}[t]{@{}l@{}}
- Maintenir les connexions temps reel avec les interfaces clientes \\
- Authentifier les connexions autorisees \\
- Diffuser les notifications de nouveaux examens/alertes \\
- Assurer la mise a jour immediate des tableaux de bord
\end{tabular} \\
\hline
\textbf{Serveur SMTP} & Acteur technique secondaire &
\begin{tabular}[t]{@{}l@{}}
- Recevoir les demandes d'envoi d'e-mails depuis le backend \\
- Envoyer les notifications pour les cas severes \\
- Fournir un canal de communication asynchrone en cas d'urgence
\end{tabular} \\
\hline
\end{tabular}
\end{table}

\textbf{Important :} \textit{Patient} et \textit{Center} ne sont pas modelises comme acteurs,
mais comme entites metier manipulees par le systeme.

\subsection{Exigences fonctionnelles}
Les exigences fonctionnelles sont regroupees par domaine :

\subsubsection*{Gestion des comptes et securite}
\begin{itemize}
    \item authentification par role (CenterAdmin, Doctor, SuperAdmin) ;
    \item gestion de session et controle d'acces selon le profil ;
    \item deconnexion securisee.
\end{itemize}

\subsubsection*{Gestion clinique et operationnelle}
\begin{itemize}
    \item creation, consultation, modification et archivage des patients ;
    \item gestion des medecins par centre ;
    \item creation d'un examen avec upload de l'image du fond d'oeil.
\end{itemize}

\subsubsection*{Analyse IA et explicabilite}
\begin{itemize}
    \item prediction automatique du grade de retinopathie (0 a 4) ;
    \item calcul d'un score de confiance associe a la prediction ;
    \item generation d'une carte de chaleur Grad-CAM ;
    \item conservation de la sortie IA dans l'historique d'examens.
\end{itemize}

\subsubsection*{Alerte et supervision}
\begin{itemize}
    \item creation d'alertes automatiques pour les cas severes (grade >= 3) ;
    \item notification temps reel du medecin ;
    \item suivi des alertes (non lue, lue, resolue) ;
    \item supervision globale et support par le SuperAdmin.
\end{itemize}

\subsection{Exigences non fonctionnelles}
Les exigences non fonctionnelles assurent la qualite globale du systeme :
\begin{itemize}
    \item \textbf{Performance} : temps de reponse court pour l'inference et l'affichage des resultats ;
    \item \textbf{Disponibilite} : acces continu a l'application pour les centres et medecins ;
    \item \textbf{Securite} : protection des donnees medicales, isolation par centre, controle des acces ;
    \item \textbf{Fiabilite} : tracabilite des examens et des alertes ;
    \item \textbf{Scalabilite} : possibilite d'ajouter des centres et d'augmenter le volume d'examens ;
    \item \textbf{Maintenabilite} : architecture modulaire (web app, service IA, service temps reel) ;
    \item \textbf{Explicabilite} : utilisation d'une visualisation interpretable pour appuyer la decision medicale.
\end{itemize}

Dans la configuration courante du projet, la visualisation explicative par defaut est une Grad-CAM
focalisee sur la couche \textit{layer4} du modele ResNet, afin de produire des cartes de chaleur plus ciblees.

\section{Modelisation des exigences fonctionnelles}

\subsection{Diagramme de cas d'utilisation general}
Le diagramme general regroupe les interactions majeures entre les trois acteurs et le systeme.

\textbf{Cas d'utilisation communs :}
\begin{itemize}
    \item se connecter ;
    \item consulter son tableau de bord ;
    \item se deconnecter.
\end{itemize}

\textbf{Cas d'utilisation CenterAdmin :}
\begin{itemize}
    \item gerer patients ;
    \item gerer medecins ;
    \item creer examen ;
    \item consulter historique des examens du centre.
\end{itemize}

\textbf{Cas d'utilisation Doctor :}
\begin{itemize}
    \item consulter liste des examens ;
    \item consulter detail d'un examen ;
    \item visualiser prediction et heatmap ;
    \item traiter alertes cliniques.
\end{itemize}

\textbf{Cas d'utilisation SuperAdmin :}
\begin{itemize}
    \item gerer centres ;
    \item gerer comptes globaux ;
    \item superviser l'activite de la plateforme ;
    \item gerer les tickets de support et le suivi global.
\end{itemize}

\subsection{Diagrammes de cas d'utilisation detaillee}

\subsubsection*{a) Use case Admin (CenterAdmin)}
Le profil CenterAdmin assure la preparation operationnelle du depistage.
Le scenario principal est le suivant :
\begin{enumerate}
    \item l'admin authentifie ouvre l'espace centre ;
    \item il cree ou selectionne un patient ;
    \item il associe un medecin referent ;
    \item il charge l'image retinienne ;
    \item il valide la creation d'examen ;
    \item le systeme declenche l'analyse IA et enregistre le resultat.
\end{enumerate}

\begin{figure}[H]
\centering
\begin{tikzpicture}[scale=0.90, every node/.style={font=\small}]
    % Limite du systeme
    \draw[thick] (1.7,0.8) rectangle (13.3,8.3);
    \node[anchor=west] at (1.9,8.02) {Systeme de depistage RD};

    % Acteur CenterAdmin
    \draw (0.1,6.0) circle (0.3);
    \draw (0.1,5.7) -- (0.1,4.8);
    \draw (-0.45,5.35) -- (0.65,5.35);
    \draw (0.1,4.8) -- (-0.35,4.1);
    \draw (0.1,4.8) -- (0.55,4.1);
    \node at (0.1,3.7) {CenterAdmin};

    % Acteur secondaire Service IA
    \draw (14.8,5.0) circle (0.3);
    \draw (14.8,4.7) -- (14.8,3.8);
    \draw (14.25,4.35) -- (15.35,4.35);
    \draw (14.8,3.8) -- (14.35,3.1);
    \draw (14.8,3.8) -- (15.25,3.1);
    \node[align=center] at (14.8,2.55) {Service IA};

    % Cas d'utilisation principaux
    \node[draw, ellipse, minimum width=4.8cm, minimum height=0.9cm] (ua1) at (7.6,7.1) {Se connecter};
    \node[draw, ellipse, minimum width=4.8cm, minimum height=0.9cm] (ua2) at (7.6,6.1) {Gerer patients};
    \node[draw, ellipse, minimum width=4.8cm, minimum height=0.9cm] (ua3) at (7.6,5.1) {Gerer medecins};
    \node[draw, ellipse, minimum width=4.8cm, minimum height=0.9cm] (ua4) at (7.6,4.1) {Creer examen};
    \node[draw, ellipse, minimum width=4.8cm, minimum height=0.9cm] (ua5) at (7.6,3.1) {Consulter historique des examens};
    \node[draw, ellipse, minimum width=4.8cm, minimum height=0.9cm] (ua6) at (7.6,2.1) {Se deconnecter};

    % Sous-cas d'utilisation de "Creer examen"
    \node[draw, ellipse, minimum width=4.2cm, minimum height=0.85cm] (ua41) at (4.2,1.25) {Televerser image retinienne};
    \node[draw, ellipse, minimum width=4.2cm, minimum height=0.85cm] (ua42) at (7.6,1.25) {Associer patient et medecin};
    \node[draw, ellipse, minimum width=4.2cm, minimum height=0.85cm] (ua43) at (11.0,1.25) {Soumettre a l'analyse IA};

    % Liens acteur principal
    \draw (0.65,5.35) -- (ua1.west);
    \draw (0.65,5.35) -- (ua2.west);
    \draw (0.65,5.35) -- (ua3.west);
    \draw (0.65,5.35) -- (ua4.west);
    \draw (0.65,5.35) -- (ua5.west);
    \draw (0.65,5.35) -- (ua6.west);

    % Include
    \draw[->, dashed] (ua4.south west) -- (ua41.north);
    \node at (5.6,2.4) {<<include>>};
    \draw[->, dashed] (ua4.south) -- (ua42.north);
    \node at (7.6,2.4) {<<include>>};
    \draw[->, dashed] (ua4.south east) -- (ua43.north);
    \node at (9.7,2.4) {<<include>>};

    % Lien avec l'acteur secondaire IA
    \draw (14.25,4.35) -- (ua43.east);
\end{tikzpicture}
\caption{Diagramme de cas d'utilisation detaille -- acteur CenterAdmin.}
\label{fig:usecase_centeradmin}
\end{figure}

\textbf{Post-condition :} un examen complet (image, grade, confiance, heatmap) est disponible.

\subsubsection*{b) Use case Doctor}
Le profil Doctor exploite les resultats pour l'aide a la decision clinique.
Le scenario principal est :
\begin{enumerate}
    \item le medecin consulte la liste des examens classes par severite ;
    \item il ouvre un examen ;
    \item il analyse le grade predit, la confiance et la heatmap ;
    \item il decide la conduite clinique (surveillance, controle rapide, orientation urgente) ;
    \item il marque l'alerte comme traitee si necessaire.
\end{enumerate}

\begin{figure}[H]
\centering
\begin{tikzpicture}[scale=0.92, every node/.style={font=\small}]
    % Limite du systeme
    \draw[thick] (1.5,0.8) rectangle (12.5,8.2);
    \node[anchor=west] at (1.7,7.95) {Systeme de depistage RD};

    % Acteur Doctor
    \draw (0,5.9) circle (0.3);
    \draw (0,5.6) -- (0,4.7);
    \draw (-0.55,5.25) -- (0.55,5.25);
    \draw (0,4.7) -- (-0.45,4.0);
    \draw (0,4.7) -- (0.45,4.0);
    \node at (0,3.6) {Doctor};

    % Cas d'utilisation
    \node[draw, ellipse, minimum width=4.9cm, minimum height=0.95cm] (uc1) at (7.0,6.9) {Se connecter};
    \node[draw, ellipse, minimum width=4.9cm, minimum height=0.95cm] (uc2) at (7.0,5.8) {Consulter liste des examens};
    \node[draw, ellipse, minimum width=4.9cm, minimum height=0.95cm] (uc3) at (7.0,4.7) {Consulter detail d'un examen};
    \node[draw, ellipse, minimum width=4.9cm, minimum height=0.95cm] (uc4) at (7.0,3.6) {Visualiser prediction et heatmap};
    \node[draw, ellipse, minimum width=4.9cm, minimum height=0.95cm] (uc5) at (7.0,2.5) {Traiter alertes cliniques};
    \node[draw, ellipse, minimum width=4.9cm, minimum height=0.95cm] (uc6) at (7.0,1.4) {Se deconnecter};

    % Associations acteur -> use cases
    \draw (0.55,5.25) -- (uc1.west);
    \draw (0.55,5.25) -- (uc2.west);
    \draw (0.55,5.25) -- (uc3.west);
    \draw (0.55,5.25) -- (uc4.west);
    \draw (0.55,5.25) -- (uc5.west);
    \draw (0.55,5.25) -- (uc6.west);

    % Relations include
    \draw[->, dashed] (uc3.south) -- (uc4.north);
    \node at (10.8,4.15) {<<include>>};
    \draw[->, dashed] (uc5.north) -- (uc3.south);
    \node at (10.8,3.05) {<<include>>};
\end{tikzpicture}
\caption{Diagramme de cas d'utilisation detaille -- acteur Doctor.}
\label{fig:usecase_doctor}
\end{figure}

\textbf{Post-condition :} la decision medicale est tracee et l'etat d'alerte est mis a jour.

\subsubsection*{c) Use case SuperAdmin}
Le profil SuperAdmin intervient au niveau de la gouvernance de la plateforme.
Le scenario principal est :
\begin{enumerate}
    \item le super administrateur se connecte au dashboard global ;
    \item il supervise les centres actifs et les comptes utilisateurs ;
    \item il controle la qualite de service (tickets, incidents, disponibilite) ;
    \item il applique les actions d'administration necessaires.
\end{enumerate}

\textbf{Post-condition :} la plateforme reste operationnelle, securisee et conforme aux objectifs de service.

\subsection{Diagramme de classes}
Le diagramme de classes represente les entites de base et leurs relations.
Les classes principales sont :
\begin{itemize}
    \item \textbf{Center} ;
    \item \textbf{User} (specialise en CenterAdmin, Doctor, SuperAdmin selon le contexte applicatif) ;
    \item \textbf{Patient} ;
    \item \textbf{Exam} ;
    \item \textbf{Alert} ;
    \item \textbf{SupportTicket} (niveau supervision).
\end{itemize}

Relations essentielles :
\begin{itemize}
    \item un Center possede plusieurs Users et plusieurs Patients ;
    \item un Patient possede plusieurs Exams ;
    \item un Exam est cree par un User et peut generer zero ou plusieurs Alerts ;
    \item une Alert est assignee a un Doctor pour traitement ;
    \item un SupportTicket est suivi par le SuperAdmin.
\end{itemize}

Cette modelisation clarifie le cycle de vie des donnees medicales, depuis l'acquisition de l'image
jusqu'au traitement des alertes.

\section{Architecture complete du systeme}

\subsection{Architecture globale}
L'architecture adoptee est une architecture modulaire multi-composants :
\begin{itemize}
    \item \textbf{Frontend Web} : interfaces de connexion, gestion metier et visualisation ;
    \item \textbf{Backend applicatif} : logique metier, API, securite, persistence ;
    \item \textbf{Service IA} : inference de grade RD et generation Grad-CAM ;
    \item \textbf{Service temps reel} : diffusion des notifications et alertes ;
    \item \textbf{Base de donnees} : stockage des utilisateurs, patients, examens et alertes.
\end{itemize}

Ce decouplage facilite la maintenance, les tests et l'evolution independante des composants.

\subsection{Architecture detaillee (diagrammes de sequence)}
Les diagrammes de sequence detaillent les echanges inter-services pour les cas critiques.

\subsubsection*{Sequence 1 : creation d'un examen}
\begin{enumerate}
    \item CenterAdmin envoie l'image au Backend ;
    \item le Backend stocke l'image et appelle le Service IA ;
    \item le Service IA retourne grade, confiance et heatmap ;
    \item le Backend enregistre l'examen ;
    \item si grade >= 3, le Backend cree une alerte ;
    \item le Service temps reel notifie le Doctor ;
    \item l'interface met a jour l'etat de maniere immediate.
\end{enumerate}

\subsubsection*{Sequence 2 : consultation et traitement d'alerte}
\begin{enumerate}
    \item le Doctor ouvre la liste des alertes ;
    \item le Backend retourne les alertes actives ;
    \item le Doctor consulte le detail d'examen ;
    \item il valide l'action clinique ;
    \item l'alerte passe a l'etat resolu avec historisation.
\end{enumerate}

\subsubsection*{Sequence 3 : supervision plateforme}
\begin{enumerate}
    \item le SuperAdmin consulte les indicateurs globaux ;
    \item le systeme agrege les statistiques multi-centres ;
    \item le SuperAdmin traite les tickets critiques et incidents ;
    \item les actions de supervision sont journalisees.
\end{enumerate}

\section{Conclusion}
Ce chapitre a presente la conception fonctionnelle et technique de la solution.
L'analyse des acteurs, la specification des exigences, les diagrammes UML et l'architecture complete
fournissent un cadre clair pour l'implementation.

La phase suivante peut s'appuyer sur cette conception pour la realisation detaillee,
les tests d'integration et la validation finale de la plateforme en contexte reel.

% =========================

\chapter{Conception}
\label{chap:conception}

\section{Introduction}
Ce chapitre presente la conception de l'application de depistage automatise de la retinopathie diabetique.
Il transforme les besoins metier en une specification technique claire, en passant par l'analyse du systeme,
la modelisation des exigences et la definition de l'architecture.

L'objectif est de garantir une solution coherent entre la pratique clinique et les contraintes informatiques,
avec une separation nette des responsabilites entre les differents profils utilisateurs.

\section{Langage et outils de modelisation}

\subsection{Langage de modelisation}
Le langage retenu pour la conception est UML (Unified Modeling Language), car il permet de representer
le systeme selon plusieurs vues complementaires :
\begin{itemize}
    \item vue fonctionnelle (diagrammes de cas d'utilisation) ;
    \item vue statique (diagramme de classes) ;
    \item vue dynamique (diagrammes de sequence).
\end{itemize}

UML est adapte a ce projet car il facilite la communication entre les acteurs techniques et metier,
et fournit une base de reference pour l'implementation.

\subsection{Outils de modelisation}
Les outils utilises pour produire et maintenir les diagrammes sont :
\begin{itemize}
    \item Mermaid (fichiers .mmd) pour versionner les diagrammes avec le code source ;
    \item draw.io / diagrams.net pour les exports visuels a integrer dans le memoire ;
    \item LaTeX pour la redaction finale et la tracabilite des sections.
\end{itemize}

Le choix de Mermaid permet une meilleure reproductibilite des diagrammes et une maintenance plus simple
au cours des iterations du projet.

\section{Analyse de systeme}

\subsection{Acteurs}
L'analyse metier identifie des acteurs principaux (metier) et des acteurs secondaires (techniques)
qui interagissent avec l'application :
\begin{table}[H]
\centering
\caption{Acteurs du systeme et principales taches}
\label{tab:acteurs_systeme}
\begin{tabular}{|p{2.8cm}|p{3.2cm}|p{8.3cm}|}
\hline
\textbf{Acteur} & \textbf{Type d'acteur} & \textbf{Taches} \\
\hline
\textbf{CenterAdmin} & Acteur metier principal &
\begin{tabular}[t]{@{}l@{}}
- Creer, consulter, modifier et archiver les fiches patient \\
- Gerer les comptes medecins rattaches au centre \\
- Verifier la qualite de l'image avant soumission \\
- Lancer l'examen et suivre son etat (en cours, termine) \\
- Consulter l'historique des examens du centre \\
- Coordonner avec le medecin en cas de resultat severe
\end{tabular} \\
\hline
\textbf{Doctor} & Acteur metier principal &
\begin{tabular}[t]{@{}l@{}}
- Consulter les examens classes par priorite clinique \\
- Analyser le grade IA, le score de confiance et la heatmap Grad-CAM \\
- Interpreter les zones suspectes avec le contexte medical du patient \\
- Valider la conduite clinique (surveillance, controle rapide, orientation) \\
- Traiter les alertes urgentes et les marquer comme resolues \\
- Ajouter des notes cliniques pour la tracabilite
\end{tabular} \\
\hline
\textbf{SuperAdmin} & Acteur de gouvernance &
\begin{tabular}[t]{@{}l@{}}
- Administrer les centres (activation, desactivation, parametrage) \\
- Gerer les comptes globaux et les droits d'acces \\
- Superviser les indicateurs globaux (volumetrie, alertes, disponibilite) \\
- Suivre les tickets de support et les incidents critiques \\
- Coordonner les actions correctives avec les equipes techniques \\
- Assurer la conformite et la securite operationnelle
\end{tabular} \\
\hline
\textbf{Service IA} & Acteur technique secondaire &
\begin{tabular}[t]{@{}l@{}}
- Recevoir l'image du fond d'oeil depuis le backend \\
- Executer l'inference pour predire le grade de retinopathie (0--4) \\
- Calculer et retourner le niveau de confiance \\
- Generer la heatmap Grad-CAM pour l'explicabilite \\
- Renvoyer une reponse structuree exploitable par l'application
\end{tabular} \\
\hline
\textbf{Service WebSocket} & Acteur technique secondaire &
\begin{tabular}[t]{@{}l@{}}
- Maintenir les connexions temps reel avec les interfaces clientes \\
- Authentifier les connexions autorisees \\
- Diffuser les notifications de nouveaux examens/alertes \\
- Assurer la mise a jour immediate des tableaux de bord
\end{tabular} \\
\hline
\textbf{Serveur SMTP} & Acteur technique secondaire &
\begin{tabular}[t]{@{}l@{}}
- Recevoir les demandes d'envoi d'e-mails depuis le backend \\
- Envoyer les notifications pour les cas severes \\
- Fournir un canal de communication asynchrone en cas d'urgence
\end{tabular} \\
\hline
\end{tabular}
\end{table}

\textbf{Important :} \textit{Patient} et \textit{Center} ne sont pas modelises comme acteurs,
mais comme entites metier manipulees par le systeme.

\subsection{Exigences fonctionnelles}
Les exigences fonctionnelles sont regroupees par domaine :

\subsubsection*{Gestion des comptes et securite}
\begin{itemize}
    \item authentification par role (CenterAdmin, Doctor, SuperAdmin) ;
    \item gestion de session et controle d'acces selon le profil ;
    \item deconnexion securisee.
\end{itemize}

\subsubsection*{Gestion clinique et operationnelle}
\begin{itemize}
    \item creation, consultation, modification et archivage des patients ;
    \item gestion des medecins par centre ;
    \item creation d'un examen avec upload de l'image du fond d'oeil.
\end{itemize}

\subsubsection*{Analyse IA et explicabilite}
\begin{itemize}
    \item prediction automatique du grade de retinopathie (0 a 4) ;
    \item calcul d'un score de confiance associe a la prediction ;
    \item generation d'une carte de chaleur Grad-CAM ;
    \item conservation de la sortie IA dans l'historique d'examens.
\end{itemize}

\subsubsection*{Alerte et supervision}
\begin{itemize}
    \item creation d'alertes automatiques pour les cas severes (grade >= 3) ;
    \item notification temps reel du medecin ;
    \item suivi des alertes (non lue, lue, resolue) ;
    \item supervision globale et support par le SuperAdmin.
\end{itemize}

\subsection{Exigences non fonctionnelles}
Les exigences non fonctionnelles assurent la qualite globale du systeme :
\begin{itemize}
    \item \textbf{Performance} : temps de reponse court pour l'inference et l'affichage des resultats ;
    \item \textbf{Disponibilite} : acces continu a l'application pour les centres et medecins ;
    \item \textbf{Securite} : protection des donnees medicales, isolation par centre, controle des acces ;
    \item \textbf{Fiabilite} : tracabilite des examens et des alertes ;
    \item \textbf{Scalabilite} : possibilite d'ajouter des centres et d'augmenter le volume d'examens ;
    \item \textbf{Maintenabilite} : architecture modulaire (web app, service IA, service temps reel) ;
    \item \textbf{Explicabilite} : utilisation d'une visualisation interpretable pour appuyer la decision medicale.
\end{itemize}

Dans la configuration courante du projet, la visualisation explicative par defaut est une Grad-CAM
focalisee sur la couche \textit{layer4} du modele ResNet, afin de produire des cartes de chaleur plus ciblees.

\section{Modelisation des exigences fonctionnelles}

\subsection{Diagramme de cas d'utilisation general}
Le diagramme general regroupe les interactions majeures entre les trois acteurs et le systeme.

\textbf{Cas d'utilisation communs :}
\begin{itemize}
    \item se connecter ;
    \item consulter son tableau de bord ;
    \item se deconnecter.
\end{itemize}

\textbf{Cas d'utilisation CenterAdmin :}
\begin{itemize}
    \item gerer patients ;
    \item gerer medecins ;
    \item creer examen ;
    \item consulter historique des examens du centre.
\end{itemize}

\textbf{Cas d'utilisation Doctor :}
\begin{itemize}
    \item consulter liste des examens ;
    \item consulter detail d'un examen ;
    \item visualiser prediction et heatmap ;
    \item traiter alertes cliniques.
\end{itemize}

\textbf{Cas d'utilisation SuperAdmin :}
\begin{itemize}
    \item gerer centres ;
    \item gerer comptes globaux ;
    \item superviser l'activite de la plateforme ;
    \item gerer les tickets de support et le suivi global.
\end{itemize}

\subsection{Diagrammes de cas d'utilisation detaillee}

\subsubsection*{a) Use case Admin (CenterAdmin)}
Le profil CenterAdmin assure la preparation operationnelle du depistage.
Le scenario principal est le suivant :
\begin{enumerate}
    \item l'admin authentifie ouvre l'espace centre ;
    \item il cree ou selectionne un patient ;
    \item il associe un medecin referent ;
    \item il charge l'image retinienne ;
    \item il valide la creation d'examen ;
    \item le systeme declenche l'analyse IA et enregistre le resultat.
\end{enumerate}

\begin{figure}[H]
\centering
\begin{tikzpicture}[scale=0.90, every node/.style={font=\small}]
    % Limite du systeme
    \draw[thick] (1.7,0.8) rectangle (13.3,8.3);
    \node[anchor=west] at (1.9,8.02) {Systeme de depistage RD};

    % Acteur CenterAdmin
    \draw (0.1,6.0) circle (0.3);
    \draw (0.1,5.7) -- (0.1,4.8);
    \draw (-0.45,5.35) -- (0.65,5.35);
    \draw (0.1,4.8) -- (-0.35,4.1);
    \draw (0.1,4.8) -- (0.55,4.1);
    \node at (0.1,3.7) {CenterAdmin};

    % Acteur secondaire Service IA
    \draw (14.8,5.0) circle (0.3);
    \draw (14.8,4.7) -- (14.8,3.8);
    \draw (14.25,4.35) -- (15.35,4.35);
    \draw (14.8,3.8) -- (14.35,3.1);
    \draw (14.8,3.8) -- (15.25,3.1);
    \node[align=center] at (14.8,2.55) {Service IA};

    % Cas d'utilisation principaux
    \node[draw, ellipse, minimum width=4.8cm, minimum height=0.9cm] (ua1) at (7.6,7.1) {Se connecter};
    \node[draw, ellipse, minimum width=4.8cm, minimum height=0.9cm] (ua2) at (7.6,6.1) {Gerer patients};
    \node[draw, ellipse, minimum width=4.8cm, minimum height=0.9cm] (ua3) at (7.6,5.1) {Gerer medecins};
    \node[draw, ellipse, minimum width=4.8cm, minimum height=0.9cm] (ua4) at (7.6,4.1) {Creer examen};
    \node[draw, ellipse, minimum width=4.8cm, minimum height=0.9cm] (ua5) at (7.6,3.1) {Consulter historique des examens};
    \node[draw, ellipse, minimum width=4.8cm, minimum height=0.9cm] (ua6) at (7.6,2.1) {Se deconnecter};

    % Sous-cas d'utilisation de "Creer examen"
    \node[draw, ellipse, minimum width=4.2cm, minimum height=0.85cm] (ua41) at (4.2,1.25) {Televerser image retinienne};
    \node[draw, ellipse, minimum width=4.2cm, minimum height=0.85cm] (ua42) at (7.6,1.25) {Associer patient et medecin};
    \node[draw, ellipse, minimum width=4.2cm, minimum height=0.85cm] (ua43) at (11.0,1.25) {Soumettre a l'analyse IA};

    % Liens acteur principal
    \draw (0.65,5.35) -- (ua1.west);
    \draw (0.65,5.35) -- (ua2.west);
    \draw (0.65,5.35) -- (ua3.west);
    \draw (0.65,5.35) -- (ua4.west);
    \draw (0.65,5.35) -- (ua5.west);
    \draw (0.65,5.35) -- (ua6.west);

    % Include
    \draw[->, dashed] (ua4.south west) -- (ua41.north);
    \node at (5.6,2.4) {<<include>>};
    \draw[->, dashed] (ua4.south) -- (ua42.north);
    \node at (7.6,2.4) {<<include>>};
    \draw[->, dashed] (ua4.south east) -- (ua43.north);
    \node at (9.7,2.4) {<<include>>};

    % Lien avec l'acteur secondaire IA
    \draw (14.25,4.35) -- (ua43.east);
\end{tikzpicture}
\caption{Diagramme de cas d'utilisation detaille -- acteur CenterAdmin.}
\label{fig:usecase_centeradmin}
\end{figure}

\textbf{Post-condition :} un examen complet (image, grade, confiance, heatmap) est disponible.

\subsubsection*{b) Use case Doctor}
Le profil Doctor exploite les resultats pour l'aide a la decision clinique.
Le scenario principal est :
\begin{enumerate}
    \item le medecin consulte la liste des examens classes par severite ;
    \item il ouvre un examen ;
    \item il analyse le grade predit, la confiance et la heatmap ;
    \item il decide la conduite clinique (surveillance, controle rapide, orientation urgente) ;
    \item il marque l'alerte comme traitee si necessaire.
\end{enumerate}

\begin{figure}[H]
\centering
\begin{tikzpicture}[scale=0.92, every node/.style={font=\small}]
    % Limite du systeme
    \draw[thick] (1.5,0.8) rectangle (12.5,8.2);
    \node[anchor=west] at (1.7,7.95) {Systeme de depistage RD};

    % Acteur Doctor
    \draw (0,5.9) circle (0.3);
    \draw (0,5.6) -- (0,4.7);
    \draw (-0.55,5.25) -- (0.55,5.25);
    \draw (0,4.7) -- (-0.45,4.0);
    \draw (0,4.7) -- (0.45,4.0);
    \node at (0,3.6) {Doctor};

    % Cas d'utilisation
    \node[draw, ellipse, minimum width=4.9cm, minimum height=0.95cm] (uc1) at (7.0,6.9) {Se connecter};
    \node[draw, ellipse, minimum width=4.9cm, minimum height=0.95cm] (uc2) at (7.0,5.8) {Consulter liste des examens};
    \node[draw, ellipse, minimum width=4.9cm, minimum height=0.95cm] (uc3) at (7.0,4.7) {Consulter detail d'un examen};
    \node[draw, ellipse, minimum width=4.9cm, minimum height=0.95cm] (uc4) at (7.0,3.6) {Visualiser prediction et heatmap};
    \node[draw, ellipse, minimum width=4.9cm, minimum height=0.95cm] (uc5) at (7.0,2.5) {Traiter alertes cliniques};
    \node[draw, ellipse, minimum width=4.9cm, minimum height=0.95cm] (uc6) at (7.0,1.4) {Se deconnecter};

    % Associations acteur -> use cases
    \draw (0.55,5.25) -- (uc1.west);
    \draw (0.55,5.25) -- (uc2.west);
    \draw (0.55,5.25) -- (uc3.west);
    \draw (0.55,5.25) -- (uc4.west);
    \draw (0.55,5.25) -- (uc5.west);
    \draw (0.55,5.25) -- (uc6.west);

    % Relations include
    \draw[->, dashed] (uc3.south) -- (uc4.north);
    \node at (10.8,4.15) {<<include>>};
    \draw[->, dashed] (uc5.north) -- (uc3.south);
    \node at (10.8,3.05) {<<include>>};
\end{tikzpicture}
\caption{Diagramme de cas d'utilisation detaille -- acteur Doctor.}
\label{fig:usecase_doctor}
\end{figure}

\textbf{Post-condition :} la decision medicale est tracee et l'etat d'alerte est mis a jour.

\subsubsection*{c) Use case SuperAdmin}
Le profil SuperAdmin intervient au niveau de la gouvernance de la plateforme.
Le scenario principal est :
\begin{enumerate}
    \item le super administrateur se connecte au dashboard global ;
    \item il supervise les centres actifs et les comptes utilisateurs ;
    \item il controle la qualite de service (tickets, incidents, disponibilite) ;
    \item il applique les actions d'administration necessaires.
\end{enumerate}

\textbf{Post-condition :} la plateforme reste operationnelle, securisee et conforme aux objectifs de service.

\subsection{Diagramme de classes}
Le diagramme de classes represente les entites de base et leurs relations.
Les classes principales sont :
\begin{itemize}
    \item \textbf{Center} ;
    \item \textbf{User} (specialise en CenterAdmin, Doctor, SuperAdmin selon le contexte applicatif) ;
    \item \textbf{Patient} ;
    \item \textbf{Exam} ;
    \item \textbf{Alert} ;
    \item \textbf{SupportTicket} (niveau supervision).
\end{itemize}

Relations essentielles :
\begin{itemize}
    \item un Center possede plusieurs Users et plusieurs Patients ;
    \item un Patient possede plusieurs Exams ;
    \item un Exam est cree par un User et peut generer zero ou plusieurs Alerts ;
    \item une Alert est assignee a un Doctor pour traitement ;
    \item un SupportTicket est suivi par le SuperAdmin.
\end{itemize}

Cette modelisation clarifie le cycle de vie des donnees medicales, depuis l'acquisition de l'image
jusqu'au traitement des alertes.

\section{Architecture complete du systeme}

\subsection{Architecture globale}
L'architecture adoptee est une architecture modulaire multi-composants :
\begin{itemize}
    \item \textbf{Frontend Web} : interfaces de connexion, gestion metier et visualisation ;
    \item \textbf{Backend applicatif} : logique metier, API, securite, persistence ;
    \item \textbf{Service IA} : inference de grade RD et generation Grad-CAM ;
    \item \textbf{Service temps reel} : diffusion des notifications et alertes ;
    \item \textbf{Base de donnees} : stockage des utilisateurs, patients, examens et alertes.
\end{itemize}

Ce decouplage facilite la maintenance, les tests et l'evolution independante des composants.

\subsection{Architecture detaillee (diagrammes de sequence)}
Les diagrammes de sequence detaillent les echanges inter-services pour les cas critiques.

\subsubsection*{Sequence 1 : creation d'un examen}
\begin{enumerate}
    \item CenterAdmin envoie l'image au Backend ;
    \item le Backend stocke l'image et appelle le Service IA ;
    \item le Service IA retourne grade, confiance et heatmap ;
    \item le Backend enregistre l'examen ;
    \item si grade >= 3, le Backend cree une alerte ;
    \item le Service temps reel notifie le Doctor ;
    \item l'interface met a jour l'etat de maniere immediate.
\end{enumerate}

\subsubsection*{Sequence 2 : consultation et traitement d'alerte}
\begin{enumerate}
    \item le Doctor ouvre la liste des alertes ;
    \item le Backend retourne les alertes actives ;
    \item le Doctor consulte le detail d'examen ;
    \item il valide l'action clinique ;
    \item l'alerte passe a l'etat resolu avec historisation.
\end{enumerate}

\subsubsection*{Sequence 3 : supervision plateforme}
\begin{enumerate}
    \item le SuperAdmin consulte les indicateurs globaux ;
    \item le systeme agrege les statistiques multi-centres ;
    \item le SuperAdmin traite les tickets critiques et incidents ;
    \item les actions de supervision sont journalisees.
\end{enumerate}

\section{Conclusion}
Ce chapitre a presente la conception fonctionnelle et technique de la solution.
L'analyse des acteurs, la specification des exigences, les diagrammes UML et l'architecture complete
fournissent un cadre clair pour l'implementation.

La phase suivante peut s'appuyer sur cette conception pour la realisation detaillee,
les tests d'integration et la validation finale de la plateforme en contexte reel.
